\documentclass[11pt]{article}

% ============================================================
%  English translation of the Italian bachelor's thesis
%  "Integrazione in termini finiti: Il Teorema di Liouville"
%  by Tommaso Mannelli Mazzoli (University of Florence, 2016-2017).
%
%  Translated into English from the Italian original by
%  Claude Opus 4.8 (an AI language model). The Italian version
%  remains the authoritative original; this translation is
%  provided for convenience and may contain errors.
% ============================================================

\usepackage[utf8]{inputenc}
\usepackage[T1]{fontenc}
\usepackage[english]{babel}
\usepackage{amsmath}
\usepackage{amssymb}
\usepackage{amsthm}
\usepackage{mathtools}
\usepackage{geometry}
\geometry{a4paper,margin=1in}

\usepackage[colorlinks=true,linkcolor=black,citecolor=black,urlcolor=blue]{hyperref}
\hypersetup{
  pdftitle={Integration in Finite Terms: Liouville's Theorem},
  pdfauthor={Tommaso Mannelli Mazzoli},
  pdfsubject={Bachelor's Thesis, University of Florence, 2016-2017},
  pdfkeywords={Translated into English from the Italian original by Claude Opus 4.8 (an AI language model). The Italian version remains the authoritative original; this translation is provided for convenience and may contain errors.}
}

% ----- theorem-like environments -----
\theoremstyle{plain}
\newtheorem{theorem}{Theorem}
\newtheorem{proposition}{Proposition}
\newtheorem{lemma}{Lemma}

\theoremstyle{definition}
\newtheorem{definition}{Definition}

\theoremstyle{remark}
\newtheorem*{remark}{Remark}
\newtheorem*{observation}{Observation}

% In the original, several statements (definitions, propositions,
% the first remarks) are introduced WITHOUT a running number.
% We reproduce those with unnumbered "starred" environments.
\newtheorem*{definition*}{Definition}
\newtheorem*{proposition*}{Proposition}

\newcommand{\C}{\mathbb{C}}
\newcommand{\R}{\mathbb{R}}
\newcommand{\Q}{\mathbb{Q}}
\newcommand{\Z}{\mathbb{Z}}
\newcommand{\N}{\mathbb{N}}
\newcommand{\Gal}{\operatorname{Gal}}
\newcommand{\Aut}{\operatorname{Aut}}
\newcommand{\Inv}{\operatorname{Inv}}
\newcommand{\Ei}{\operatorname{Ei}}
\newcommand{\Li}{\operatorname{Li}}
\renewcommand{\deg}{\operatorname{deg}}
\DeclareMathOperator{\Real}{Re}
\DeclareMathOperator{\Imag}{Im}

\begin{document}

% ============================================================
%                       TITLE PAGE
% ============================================================
\begin{titlepage}
\centering

{\Large\bfseries Universit\`a degli Studi di Firenze}\\[2pt]
{\small (University of Florence)}\\[4pt]
\rule{\textwidth}{0.4pt}\\[6pt]

{\sc School of Mathematical, Physical and Natural Sciences}\\[2pt]
{Bachelor's Degree Programme in Mathematics}\\[40pt]

{\scshape Bachelor's Thesis}\\[60pt]

{\LARGE\bfseries Integration in Finite Terms:\\[6pt]
Liouville's Theorem.}\\[16pt]

{\large\itshape Integrazione in termini finiti:\\
Il Teorema di Liouville.}\\[10pt]
{\small (Italian original title)}\\[70pt]

\begin{flushleft}
\hspace{0.04\textwidth}%
\begin{minipage}[t]{0.45\textwidth}
\textbf{Candidate:}\\[4pt]
\textbf{Tommaso Mannelli Mazzoli}\\[2pt]
\textbf{\small Student no.\ 5650541}
\end{minipage}%
\hfill
\begin{minipage}[t]{0.45\textwidth}
\textbf{Advisor (Relatore):}\\[4pt]
\textbf{Prof.\ Andrea Colesanti}
\end{minipage}
\end{flushleft}

\vfill

{\bfseries Academic Year 2016--2017}\\[14pt]

\rule{\textwidth}{0.4pt}\\[6pt]
\begin{minipage}{0.92\textwidth}
\footnotesize\itshape
Translated into English from the Italian original by Claude Opus 4.8
(an AI language model). The Italian version remains the authoritative
original; this translation is provided for convenience and may contain
errors.
\end{minipage}

\end{titlepage}

% ============================================================
%                    TABLE OF CONTENTS
% ============================================================
\renewcommand{\contentsname}{Contents}
\tableofcontents
\newpage

% ============================================================
\section{Introduction}
% ============================================================
Every student of Calculus~I who tries to integrate functions such as
$e^{x^2}$, $\sin(x^2)$, $e^{1/x}$ is told that they are not
elementarily integrable, that is, it is not possible to write their
antiderivative by combining in a finite way the four elementary
operations, raising to powers and extracting roots, logarithms,
exponentials, trigonometric and inverse trigonometric functions. But
why? Let us give another example of an integrable function that does
not admit an elementary antiderivative. In number theory there is much
interest in the study of the function
$\pi(x) = |\{1 \le p \le x : p \text{ prime}\}|$ of the real variable
$x$, which counts the number of primes up to $x$. Defining the
logarithmic integral function $\Li(x) = \int_{2}^{x} \frac{dt}{\ln(t)}$,
it has been proved that, as $x \to \infty$, $\pi(x) \sim \Li(x)$. This
example shows how important the computation of $\Li(x)$ can be, which
however is not expressible in terms of elementary functions. How is it
possible to prove a result like this? Certainly, first of all it is
necessary to give a precise definition of an elementary function, and
the tool that will allow us to move within this domain with ease is
differential algebra. Thanks to it we will prove results (first and
foremost Liouville's theorem) with which we will give a
characterisation of the elementarily integrable functions.

% ============================================================
\section{Prerequisites}
% ============================================================
The notions reported in this section are taken from~\cite{Casolo}.

A \emph{field} $F$ is a commutative ring with unit in which every
nonzero element is invertible; we denote by $F^{*} = F \setminus \{0\}$.
Let $F, E$ be fields. $E$ is said to be an \emph{extension} of $F$
(and one writes $E\mid F$) if there exists an injective field
homomorphism (called an \emph{embedding}):
\[
\sigma : F \longrightarrow E.
\]
Let $F$ be a field; we define by $F[x]$ the ring of polynomials with
coefficients in $F$ and by $F(x)$ the field of fractions of $F[x]$.
Thus
\[
F(x) = \left\{ \frac{f}{g} : f, g \in F[x],\ g \neq 0 \right\}.
\]
Note that $F(x)$ is a field and $F(x)\mid F$ is a field extension.

\begin{definition*}
Let $E\mid F$ be a field extension. Let $b \in E$. We say that $b$ is
\textbf{algebraic} over $F$ if there exists a polynomial $f \neq 0$ in
$F[x]$ such that $f(b) = 0$. An extension $E\mid F$ is said to be an
\textbf{algebraic extension} if every element of $E$ is algebraic over
$F$.
\end{definition*}

\begin{definition*}
If $E\mid F$ is a field extension, then it is possible to regard $E$ in
a natural way as a vector space over $F$. The dimension of $E$ as a
vector space over $F$ is called the \textbf{degree} of $E$ over $F$ and
is denoted by $[E : F]$. If $[E : F] < \infty$, then $E\mid F$ is said
to be a \textbf{finite} extension.
\end{definition*}

We also observe that every algebraic extension is finite.

\begin{definition*}
A field extension $E\mid F$ is said to be \textbf{normal} if every
element $b \in E$ is algebraic over $F$ and its minimal polynomial
factors completely in $E[x]$.
\end{definition*}

\begin{definition*}
A polynomial $f$ is said to be \emph{separable} if each of its
irreducible factors has all its roots simple (in an extension
containing all its roots). A field extension $E\mid F$ is said to be
\textbf{separable} (over $F$) if it is algebraic and for every
$b \in E$, the minimal polynomial $\min_{F}(b)$ is separable over $F$.
\end{definition*}

\begin{definition*}
Let $R$ be a ring and let $1 \in R$ be its multiplicative unit. The
\emph{characteristic} of $R$ is $0$ if $n \cdot 1 \neq m \cdot 1$ for
all $n, m \in \Z$; otherwise it is the smallest natural number $n > 0$
such that $n \cdot 1 = 0$.
\end{definition*}

\begin{proposition*}
Let $F$ be a field of characteristic $0$, or a finite field. Then
every extension $E\mid F$ is separable.
\end{proposition*}

Every field that we consider from now on will have characteristic $0$
(as for example $\Q$, $\R$ or $\C$).

\begin{definition*}
A field extension $E\mid F$ which is finite, normal and separable is
said to be a \textbf{Galois} extension.
\end{definition*}

Let us give some reminders about Galois theory, which will be used in
what follows.

\begin{definition*}
Let $E\mid F$ be a field extension. The set of the $F$-automorphisms of
$E$ (that is, of the automorphisms of $E$ that fix every element of $F$)
is a subgroup of the group $\Aut(E)$ and is denoted by $\Gal(E\mid F)$.
Moreover, given $f \in F[x]$, one defines the \emph{Galois group} of
$f$ over $F$ as $\Gal(E\mid F)$ with $E$, where $E$ is the splitting
field of $f$ over $F$. Now let $H$ be a subgroup of $\Gal(E\mid F)$; we
define the \emph{field of invariants}
$\Inv_{E}(H) = \{ b \in E \mid \sigma(b) = b\ \forall \sigma \in
\Gal(E\mid F) \}$.
\end{definition*}

For the proof of the following two propositions see~\cite{Casolo}.

\begin{proposition}\label{prop:transitive}
Let $E$ be the splitting field of a separable polynomial
$f \in F[x]$ of degree $n$. Then the subgroup of $S_{n}$ isomorphic to
the Galois group of $f$ over $F$ is transitive if and only if $f$ is
irreducible over $F$.
\end{proposition}

\begin{proposition}\label{prop:inv}
Let $E\mid F$ be a Galois extension; then $\Inv_{E}(E) = F$.
\end{proposition}

\subsection*{Symmetric polynomials in $F[x_{1}, \ldots, x_{n}]$}

\begin{definition*}
Let $F$ be a field, $f \in F[x_{1}, \ldots, x_{n}]$. We say that $f$ is
\emph{symmetric} if for every permutation $\sigma \in S_{n}$ (the group
of permutations of $n$ elements) one has that
$f(x_{\sigma(1)}, \ldots, x_{\sigma(n)}) = f(x_{1}, \ldots, x_{n})$.
\end{definition*}

\begin{theorem}\label{thm:sym}
Let $F$ be a field, $f \in F[x]$ a monic polynomial of degree $n > 0$
with roots $a_{1}, \ldots, a_{n}$ in one of its splitting fields over
$F$. Then for every symmetric polynomial
$p \in F[x_{1}, \ldots, x_{n}]$ one has that
$p(a_{1}, \ldots, a_{n}) \in F$.
\end{theorem}

\begin{proof}
Let $f$ be a polynomial as in the hypotheses and let $E$ be its
splitting field; then $E\mid F$ is a Galois extension. We define
$G := \Gal(E\mid F)$ and we consider $p \in F[x_{1}, \ldots, x_{n}]$ a
symmetric polynomial. Every element permutes the roots of $f$, hence
for every $\sigma \in \Gal(E\mid F)$ one has that
$\sigma(p(a_{1}, \ldots, a_{n})) = p(a_{\sigma(1)}, \ldots,
a_{\sigma(n)}) = p(a_{1}, \ldots, a_{n})$, and therefore, by
Proposition~\ref{prop:inv}, it follows that
$p \in \Inv_{G}(G) = F$.
\end{proof}

\subsection*{Partial fraction decomposition}
A standard procedure for the computation of the integrals of rational
functions of the type
\[
\int \frac{dx}{x^{2}-1}
\]
is to break up the integrand (which we do not know how to integrate)
into a sum of other, simpler rational functions (which we do know how
to integrate). For example, in our case
\[
\int \frac{dx}{x^{2}-1}
= \int \frac{1}{2}\left( \frac{1}{x-1} - \frac{1}{x+1} \right) dx
= \int \frac{1}{2(x-1)}\, dx - \int \frac{1}{2(x+1)}\, dx.
\]
But is it always possible to act in this way no matter which rational
function is chosen, that is, for any polynomials $P, Q$ in $F[x]$?

Let us begin by fixing a field $F$ and choosing a rational function
$v = \frac{P}{Q}$ with $P, Q \in F[x]$ coprime. We factor $Q$ into
irreducible factors:
\[
Q = Q_{1}^{k_{1}} \cdot Q_{2}^{k_{2}} \cdots Q_{N}^{k_{N}}.
\]
Note in particular that $Q_{1}^{k_{1}}$ and
$Q_{2}^{k_{2}} \cdots Q_{N}^{k_{n}}$ are coprime, and therefore by
B\'ezout's identity we know that there exist two polynomials $T$ and
$S$ such that $T Q_{1}^{k_{1}} + S Q_{2}^{k_{2}} \cdots Q_{N}^{k_{N}} =
1$, and hence we can write
\[
v = \frac{P}{Q}
= \frac{P\bigl(T Q_{1}^{k_{1}} + S Q_{2}^{k_{2}} \cdots
Q_{N}^{k_{N}}\bigr)}{Q_{1}^{k_{1}} \cdot Q_{2}^{k_{2}} \cdots
Q_{N}^{k_{N}}}
= \frac{P T}{Q_{2}^{k_{2}} \cdots Q_{N}^{k_{N}}}
+ \frac{P S}{Q_{1}^{k_{1}}}.
\]
What have we gained? We have obtained two rational functions in which
one denominator is a power of an irreducible polynomial, while the
other denominator has one factor fewer than $Q$ in its factorisation.
Applying this procedure $N-1$ times we obtain an expression of the type
\[
v = \frac{T_{1}}{Q_{1}^{k_{1}}} + \frac{T_{2}}{Q_{2}^{k_{2}}}
+ \cdots + \frac{T_{N}}{Q_{N}^{k_{N}}}
= \sum_{i=1}^{N} \frac{T_{j}}{Q_{j}^{k_{j}}}
\]
where each $Q_{j}$ is irreducible. If now the degree of $T_{j}$ were
greater than or equal to the degree of $Q_{j}$, we could use Euclid's
algorithm to write $T_{j} = U_{j} Q_{j} + R_{j}$ where $R_{j}$ is a
polynomial of degree less than that of $Q_{j}$. As a result we obtain
\[
\frac{T_{j}}{Q_{j}^{k_{j}}}
= \frac{U_{j}}{Q_{j}^{k_{j}-1}} + \frac{R_{j}}{Q_{j}^{k_{j}}}.
\]
In a finite number of steps we arrive at an identity of the form
\[
\frac{R_{j}}{Q_{j}^{k_{j}}}
= S_{0,j} + \sum_{l=1}^{m_{j}} \frac{S_{l,j}}{Q_{j}^{l}}
\]
where, for $l \ge 1$, each polynomial $S_{l,j}$ has degree strictly
less than $Q_{j}$. We have thus proved the following theorem:

\begin{theorem}\label{thm:partialfractions}
Let $F$ be a field, $v \in F(x)$ and $P, Q \in F[x]$ such that
$v = \frac{P}{Q}$. Let $Q = Q_{1}^{k_{1}} \cdots Q_{N}^{k_{N}}$ be its
factorisation into irreducibles. Then there exist polynomials
$T, S_{ij} \in F[x]$ such that the degree of $S_{i,j}$ is strictly less
than the degree of $Q_{j}$ for every $i \ge 1$, and moreover
\[
v = T + \sum_{i=1}^{n} \sum_{j=1}^{k_{i}} \frac{S_{ij}}{Q_{i}^{j}}.
\]
Such an expression of $v$ is called the \emph{partial fraction
decomposition}.
\end{theorem}

% ============================================================
\section{Differential Fields}
% ============================================================
\subsection*{Definitions and properties of differential fields}

\begin{definition}\label{def:diffring}
A \textbf{differential ring} is a pair $(A, ')$ where $A$ is a ring and
$' : A \to A$ a map (called a \emph{derivation}) such that:
\begin{itemize}
\item[(D1)] $(x + y)' = x' + y'$ for all $a, b \in A$;
\item[(D2)] $(xy)' = x' y + x y'$ for all $a, b \in A$.
\end{itemize}
If $A$ is a field, we will say that $(A, ')$ is a \textbf{differential
field}.
\end{definition}

In general a ring can have more than one derivation defined on it. For
example, the ring of polynomials with rational coefficients in two
indeterminates $\Q[x, y]$ has at least three derivations: $0$, $d/dx$,
$d/dy$; but in fact it has many more: every linear combination of
derivations with coefficients in the ring is a derivation on the ring.
An example of a differential field, which is then the one that will
accompany us along this path, is $(\C(x), ')$ where $'$ is the usual
derivative. Indeed $\C(x)$ is a field and $'$ a map that satisfies
properties (D1) and (D2).

\begin{proposition}\label{prop:derivrules}
Let $(F, ')$ be a differential field; then
\begin{enumerate}
\item $1_{F}' = 0_{F}' = 0_{F}$;
\item $(x^{k})' = k x^{k-1} x'$ for every $x \in F$ and every
$k \in \Z$;
\item $(-x)' = -x'$ for every $x \in F$;
\item $\left( \dfrac{x}{y} \right)' = \dfrac{x' y - x y'}{y^{2}}$ for
every $x, y \in F$, $y \neq 0$;
\item $\dfrac{\bigl(x_{1}^{k_{1}} x_{2}^{k_{2}} \cdots
x_{n}^{k_{n}}\bigr)'}{x_{1}^{k_{1}} x_{2}^{k_{2}} \cdots x_{n}^{k_{n}}}
= k_{1} \dfrac{x_{1}'}{x_{1}} + k_{2} \dfrac{x_{2}'}{x_{2}}
+ \ldots k_{n} \dfrac{x_{n}'}{x_{n}}$ for every
$x_{1}, \ldots x_{n} \in F^{*}$ and $k_{1}, \ldots, k_{n} \in \Z$.
\end{enumerate}
\end{proposition}

A (rather routine) proof of these properties can be found
in~\cite{DeLellis}.

\begin{definition}\label{def:constants}
We call the \textbf{subfield of constants of} $(F, ')$ the set
\[
C_{F} = \{ x \in F : x' = 0 \}.
\]
\end{definition}

\begin{observation}
It is not difficult to prove that $C_{F}$ is a field.
\end{observation}

\subsection*{Extensions of differential fields}

\begin{definition}\label{def:diffext}
Let $(F, ')$ and $(E, {}^{*})$ be two differential fields. We say that
$(E, {}^{*})$ is a \textbf{differential extension} of $(F, ')$ if the
following properties hold:
\begin{itemize}
\item[(i)] $E$ is an extension of $F$. In symbols $E\mid F$;
\item[(ii)] ${}^{*}$ extends $'$. In symbols $x^{*} = x'$ for every
$x \in F$.
\end{itemize}
\end{definition}

Let $(F, ')$ be a differential field. We consider the polynomial ring
$F[x]$ and the maps $D_{0}, D_{1} : F[x] \to F[x]$ defined as follows:
\[
D_{0}\left( \sum_{i=0}^{n} a_{i} x^{i} \right)
= \sum_{i=0}^{n} a_{i}' x^{i};
\qquad
D_{1}\left( \sum_{i=0}^{n} a_{i} x^{i} \right)
= \sum_{i=1}^{n} i a_{i} x^{i-1}.
\]
It is not difficult to prove the following proposition:

\begin{proposition}\label{prop:D0D1}
$D_{0}$ and $D_{1}$ are derivations in the ring $F[x]$.
\end{proposition}

The following theorem shows how, given an algebraic extension, it is
possible to extend the derivation on the base field in a unique manner.

\begin{theorem}\label{thm:extend}
Let $(F, ')$ be a differential field, let $E\mid F$ be an algebraic
extension. Then there exists a unique map ${}^{*} : E \longrightarrow E$
such that $(E, {}^{*})$ is a differential extension of $(F, ')$.
\end{theorem}

\begin{proof}
We prove only uniqueness (the proof of existence can be found
in~\cite{Bronstein}). Suppose that there exists a derivation ${}^{*}$
in $E$ that extends the one on $F$. Let $b \in E$ and let $f \in F[x]$
be its minimal polynomial over $F$ (it exists since $b$ is algebraic
over $F$). Then, using the properties of the derivation, one has that
\[
0 = (f(b))^{*}
= \left( \sum_{i=0}^{n} a_{i} b^{i} \right)^{*}
= \sum_{i=0}^{n} a_{i}' b^{i}
+ b^{*} \sum_{i=1}^{n} a_{i} b^{i-1} i
= (D_{0} f)(b) + b^{*} (D_{1} f)(b),
\]
from which one obtains that
\[
b^{*} = - \frac{(D_{0} f)(b)}{(D_{1} f)(b)}.
\]
(Note that, since $f$ is the minimal polynomial of $b$ over $F$, then
$D_{1} f$ is a polynomial of degree less than $f$, hence it cannot be
zero.) Thus, if there exists a derivation extending $'$ on $F[b]$, it
is unique. Now, since $E\mid F$ is algebraic, for every $b \in E$ one
has that the derivation extends to $F[b]$ in a unique manner. By the
arbitrariness of $b$ one obtains the thesis.
\end{proof}

\begin{observation}
The previous theorem also provides us with a way to compute the
derivative of algebraic elements in a purely abstract manner. Let us
think for example of the differential field of fractions of polynomials
with coefficients in $\C$ with the usual derivative $(\C(x), ')$. The
element $b = \sqrt{x}$ does not belong to the field but is a root of
the polynomial $P(t) = t^{2} - x \in \C(x)[t]$. We note moreover that,
since $P$ is monic and irreducible, it is also the minimal polynomial
of $b$ over $\C(x)$. One immediately finds that $(D_{0})(b) = -1$ and
$(D_{1})(b) = 2b$. Hence, by what was proved above,
\[
(\sqrt{x})' = b^{*} = -\frac{-1}{2b} = \frac{1}{2\sqrt{x}}.
\]
\end{observation}

% ============================================================
\section{Liouville's Theorem}
% ============================================================

\begin{definition}\label{def:logexp}
Let $(F, ')$ be a differential field and $E$ a differential extension
of it. An element $t \in E$ is:
\begin{itemize}
\item a \textbf{logarithm} over $F$ if there exists $b \in F^{*}$ such
that $t' = b'/b$, and we write $t = \log(b)$;
\item an \textbf{exponential} over $F$ if there exists $b \in F$ such
that $t' = t b'$ and we write $t = e^{b}$;
\item \textbf{elementary} over $F$ if it is algebraic or an exponential
or a logarithm.
\end{itemize}
\end{definition}

\begin{definition}\label{def:elemext}
$E$ is an \textbf{elementary extension} of $F$ if there exist
$t_{1}, \ldots t_{n}$ in $E$ such that $E = F(t_{1}, \ldots, t_{n})$ and
one has that $t_{i}$ is elementary over
$F(t_{1}, t_{2}, \ldots, t_{i-1})$. We will say moreover that
$f \in F$ is \textbf{elementarily integrable} over $F$ if there exists
an elementary extension $E\mid F$ and $g \in E$ such that $g' = f$.
Finally, an \textbf{elementary function} is any element of an
elementary extension of $(\C(x), d/dx)$.
\end{definition}

For example $e^{(e^{x})}$ is elementary (two exponential extensions of
$\C(x)$), as is $\sqrt{\log(x)}$ (a logarithmic extension followed by
an algebraic one).

We can now define precisely the problem of integration in finite
terms: given a differential field $F$ and $f \in F$, the problem
consists in deciding in a finite number of steps whether $f$ is
elementarily integrable over $F$ and in computing its antiderivative,
if it exists.

\begin{observation}
Consider $F = \C(x, t_{1}, t_{2})$ where $x, t_{1}, t_{2}$ are
indeterminates over $\C$, with the derivation $D$ such that $Dx = 1$,
$Dt_{1} = t_{1}$, $Dt_{2} = t_{1}/x$ (in practice $t_{1} = e^{x}$ and
$t_{2} = \Ei(x)$ with $\Ei(x) := \int e^{x}/x\, dx$, which we will
prove later not to be elementarily integrable). Then
\[
\int \frac{t_{1} t_{2}}{x}\, dx = \frac{t_{2}^{2}}{2} \in F,
\]
and hence the element $\frac{t_{1} t_{2}}{x}$ (which morally is
$\frac{e^{x} \Ei(x)}{x}$) is elementarily integrable over $F$ (after
all, its integral lies in $F$, which is certainly an elementary
extension of $F$) but is not an elementary function, since it is not an
element of an elementary extension of $(\C(x), d/dx)$.
\end{observation}

\begin{observation}
The elementary functions of Definition~\ref{def:elemext} include all
the usual elementary functions of analysis, since the trigonometric
functions and their inverses can be rewritten in terms of complex
exponentials and logarithms (thanks to the celebrated formula of Euler
$e^{fi} = \cos(f) + i \sin(f)$).
\end{observation}

\begin{lemma}\label{lem:degree}
Let $(F, ')$ be a differential field. Let $F(t)$ be a differential
extension of $F$ with $t$ transcendental over $F$ such that
$C_{F} = C_{F(t)}$.
\begin{enumerate}
\item If $t' \in F$, then for every polynomial $f(t) \in F[t]$ of
degree $n > 0$, $(f(t))'$ is a polynomial in $F[t]$ of degree $n$ or
$n-1$.
\item If $t'/t \in F$, then for every $a \in F$ one has that
$(a t^{n}) = h t^{n}$ for some $h \in F$, and for every polynomial
$f(t) \in F[t]$ one has that $(f(t))'$ is a polynomial in $F[t]$ of the
same degree as $f(t)$.
\end{enumerate}
\end{lemma}

\begin{proof}
1. Let $t' = b \in F$ and $n = \deg f(t)$. Then we can write
$f(t) = a_{n} t^{n} + a_{n-1} t^{n-1} + \cdots + a_{0}$ with
$a_{i} \in F$ and $a_{n} \neq 0$. Then
\[
f'(t) = a_{n}' t^{n} + (a_{n} b n + a_{n-1}') t^{n-1} + \cdots
+ (a_{n-i} b (n-i) + a_{n-i-1}') t^{n-i} + \cdots + a_{1}'.
\]
Hence the derivative of $f(t)$ continues to be a polynomial in $F[t]$
of degree $n$ if $a_{n}' \neq 0$. Suppose now that $a_{n}' = 0$ and
suppose by contradiction that $(a_{n} b n + a_{n-1}') = 0$; one would
then have that $a_{n} b n + a_{n-1} = (n a_{n} t + a_{n-1})' = 0$, and
hence $n a_{n} t + a_{n-1}$ would be a constant, that is, an element of
$C_{F(t)}$. But since we have assumed $C_{F} = C_{F(t)}$, one would
have that $n a_{n} t + a_{n-1} \in F$, that is, $t \in F$, which is
absurd since $t$ is transcendental over $F$; hence if $a_{n}$ is a
constant, then the degree of $(f(t))'$ is $n-1$.

2. Suppose now that $t'/t = b \in F$. Let $0 \neq a \in F$. Then
\[
(a t^{n})' = a' t^{n} + a n t^{n-1} t' = (a' + a b n) t^{n}.
\]
Now if $a' + a b n = 0$, then $(a t^{n})' = 0$, and hence $a t^{n}$
would be a constant and therefore an element of $F$, contradicting the
hypothesis of transcendence of $t$ over $F$. Hence $a' + a b n \neq 0$.
Therefore one has that if $f(t) \in F[t]$, then $(f(t))'$ is a
polynomial of the same degree as $f(t)$.
\end{proof}

We now prove the central result of the thesis.

\begin{theorem}[Liouville]\label{thm:liouville}
Let $(F, ')$ be a differential field and let $f \in F$. The following
two assertions are equivalent:
\begin{itemize}
\item[(i)] There exists an elementary extension $E\mid F$ such that
$C_{E} = C_{F}$ and $g \in E$ with $g' = f$.
\item[(ii)] There exist $v \in F$, $u_{1}, \ldots, u_{n} \in F^{*}$ and
$c_{1}, \ldots, c_{n} \in C_{F}$ such that
\[
f = v' + \sum_{i=1}^{n} c_{i} \frac{u_{i}'}{u_{i}}.
\]
\end{itemize}
\end{theorem}

\begin{proof}
(ii) $\Rightarrow$ (i). It suffices to set
\[
g = v + \sum_{i=1}^{N} c_{i} \log(u_{i}).
\]
In this way $g' = f$.

(i) $\Rightarrow$ (ii). Since the extension $E\mid F$ is elementary,
there exist $t_{1}, t_{2}, \ldots, t_{m} \in F$ such that
$E = F(t_{1}, \ldots, t_{m})$ and, for every $i = 1, \ldots, m$, $t_{i}$
is an algebraic, or exponential, or logarithmic element over
$F(t_{1}, \ldots, t_{i-1})$. Moreover, since $E$ has the same subfield
of constants as $F$, then every differential field $L$ such that
$E\mid L\mid F$ will have the same subfield of constants. We prove the
theorem by induction on $m$. For $m = 0$ one has that $E = F$ and hence
the theorem is trivially proved: with $v = y$. We assume the theorem
true for every elementary field extension $F(t_{1}, \ldots, t_{m})$ and
we consider the chain of extensions
$F(t_{1}, \ldots, t_{m}) \mid F(t_{1}) \mid F$. By the inductive
hypothesis there exist $c_{1}, \ldots, c_{n} \in C_{F}$,
$u_{1}, \ldots, u_{n} \in F(t_{1})$, $v \in F(t_{1})$ such that
$f = v' + \sum_{i=1}^{n} c_{i} \frac{u_{i}'}{u_{i}}$. We must prove that
such an expression for $f$ exists with $u_{1}, \ldots, u_{n} \in F$,
$v \in F$ (possibly for a different $n$). Let $t = t_{1}$; we divide the
proof into three cases: the one in which $t$ is algebraic, the one in
which $t$ is a logarithm, and finally the one in which $t$ is an
exponential.

\medskip
\noindent\textbf{The algebraic case.}
Since $t$ is algebraic over $F$, there exist
$U_{1}, \ldots, U_{n}, V \in F[x]$ such that $U_{1}(t) = u_{1}, \ldots,
U_{n}(t) = u_{n}$, $V(t) = v$. We now consider the minimal polynomial
of $t$ over $F$, $p_{t} = \min_{F}(b)$, and the distinct conjugates of
$t$, namely $\tau_{1}(= t), \tau_{2}, \ldots, \tau_{s}$ in $K$, the
splitting field of $p_{t}$ over $F$. Since $K\mid F$ is algebraic, by
Theorem~\ref{thm:extend} the derivation on $F$ can be extended in a
unique manner on $F$. By hypothesis one then has that
\begin{equation}\label{eq:1}
f = (V(\tau_{1}))' + \sum_{i=1}^{n} c_{i}
\frac{(U_{i}(\tau_{1}))'}{U_{i}(\tau_{1})}.
\end{equation}
Since the polynomial $p_{t}$ is separable and irreducible over $F$,
then by Proposition~\ref{prop:transitive} the subgroup of $S_{n}$
isomorphic to the Galois group $\Gal(K\mid F)$ is transitive, and hence
for every $\tau_{j}$ conjugate of $\tau_{1}$ there exists
$\sigma \in \Gal(K\mid F)$ such that $\sigma(\tau_{1}) = \tau_{j}$.
Applying $\sigma$ to both members of equation~\eqref{eq:1} one has:
\begin{align*}
\sigma(f) = f
&= \sigma\!\left( (V(\tau_{1}))' + \sum_{i=1}^{n} c_{i}
\frac{(U_{i}(\tau_{1}))'}{U_{i}(\tau_{1})} \right)\\
&= \sigma\bigl((V(\tau_{1}))'\bigr) + \sum_{i=1}^{n} \sigma(c_{i})
\frac{\sigma(U_{i}(\tau_{1}))'}{\sigma(U_{i}(\tau_{1}))}\\
&= \sigma(V(\tau_{1}))' + \sum_{i=1}^{n} c_{i}
\frac{\sigma((U_{i}(\tau_{1}))'}{U_{i}(\tau_{j})}.
\end{align*}
With the notation of Theorem~\ref{thm:extend} we write more explicitly
what $\sigma(V(\tau_{1})')$ is:
\begin{align*}
\sigma(V(\tau_{1})')
&= \sigma\bigl((D_{0} V)(\tau_{1}) + (D_{1} V)(\tau_{1}) \tau_{1}'\bigr)\\
&= D_{0} V(\sigma(\tau_{1})) + D_{1} V(\sigma(\tau_{1}))
\left( -\frac{(D_{0} V)(p_{t})}{(D_{1} V)(\sigma(p_{t}))} \right)\\
&= D_{0} V(\tau_{j}) + D_{1} V(\tau_{j})
\left( -\frac{D_{0} V(p_{t})}{D_{1} V(p_{t})} \right)
= (V(\tau_{j}))'.
\end{align*}
In an analogous way one has that
$\sigma(U_{i}(\tau_{1}))' = (U_{i}(\tau_{j}))'$. One thus has that
\[
f = (V(\tau_{j}))' + \sum_{i=1}^{n} c_{i}
\frac{(U_{i}(\tau_{j}))'}{U_{i}(\tau_{j})},
\qquad \forall\, j \in \{1, \ldots, s\}.
\]
Now summing over $j$ on both members and dividing by $s$, one obtains
\begin{align*}
f &= \frac{(V(\tau_{1}))'}{s} + \cdots + \frac{(V(\tau_{s}))'}{s}
+ \sum_{i=1}^{n} \frac{c_{i}}{s}
\left( \frac{(U_{i}(\tau_{1}))'}{U_{i}(\tau_{1})} + \cdots
+ \frac{(U_{i}(\tau_{s}))'}{U_{i}(\tau_{s})} \right)\\
&= \left( \frac{V(\tau_{1}) + \cdots + V(\tau_{s})}{s} \right)'
+ \sum_{i=1}^{n} \frac{c_{i}}{s}
\frac{(U_{i}(\tau_{1}) \cdots U_{i}(\tau_{s}))'}{U_{i}(\tau_{1})
\cdots U_{i}(\tau_{s})},
\end{align*}
where in the second passage point~5 of
Proposition~\ref{prop:derivrules} was used. Now we note that
$U_{i}(x_{1}) \cdots U_{i}(x_{s})$ and $V(x_{1}) + \cdots + V(x_{s})$
are symmetric polynomials in $F[x_{1}, \ldots, x_{s}]$, and since
$\tau_{1} \ldots, \tau_{s}$ are the roots of the minimal polynomial of
$t$ over $F$, then by Theorem~\ref{thm:sym} one has that
$U_{i}(\tau_{1}) \ldots U_{i}(\tau_{s})$,
$V(\tau_{1}) + \cdots + V(\tau_{s}) \in F$, and hence one has the
thesis.

\medskip
\noindent
Before passing to the logarithm and exponential cases, it is convenient
to make a preliminary observation. Suppose that $t$ is transcendental
over $F$. Then by the inductive hypothesis we know that there exist
$u_{1}(t), \ldots, u_{n}(t), v(t) \in F(t)$ and
$c_{1}, \ldots, c_{n} \in C_{F}$ such that
\begin{equation}\label{eq:2}
f = (v(t))' + \sum_{i=1}^{n} c_{i} \frac{(u_{i}(t))'}{u_{i}(t)}.
\end{equation}
Now, since $u_{i}(t) \in F(t)$ for every $i = 1, \ldots, n$, there
exist $p_{i}, q_{i} \in F[t]$ such that $u_{i}(t) = \frac{p_{i}}{q_{i}}$.
Considering the factorisation into irreducibles of $p_{i}$ and $q_{i}$
one obtains
\[
u_{i}(t) = \frac{p_{i}}{q_{i}}
= \frac{p_{1i}^{\alpha_{1i}} \cdots p_{si}^{\alpha_{si}}}{q_{1i}^{\beta_{1i}}
\cdots q_{ri}^{\beta_{ri}}}
= p_{1i}^{\alpha_{1i}} \cdots p_{si}^{\alpha_{si}}
q_{1i}^{-\beta_{1i}} \cdots q_{ri}^{-\beta_{ri}},
\]
from which it follows, using the properties of
Proposition~\ref{prop:derivrules}, that
\[
\frac{(u_{i}(t))'}{u_{i}(t)}
= \sum_{j=1}^{s} \alpha_{ji} \frac{p_{ji}'}{p_{ji}}
- \sum_{k=1}^{r} \beta_{ki} \frac{q_{ki}'}{q_{ri}}.
\]
From this observation it follows that the summation in
equation~\eqref{eq:2} can be rewritten in the form
\[
\sum_{i=1}^{N} \tilde{c}_{i} \frac{\tilde{u}_{i}'}{\tilde{u}_{i}},
\]
in which for every $i = 1, \ldots N$, $\tilde{c}_{i} \in C_{F}$ and
$\tilde{u}_{i} \in F[t]$ is monic and irreducible (or directly an
element of $F$). We now consider Theorem~\ref{thm:partialfractions},
which tells us that it is possible to decompose $v$ into partial
fractions:
\[
v = T + \sum_{i=1}^{n} \sum_{j=1}^{k_{i}} \frac{S_{ij}}{P_{i}^{j}}.
\]
We can therefore rewrite equation~\eqref{eq:2} in the following way:
\begin{equation}\label{eq:3}
f = T' + \sum_{i=1}^{N} \tilde{c}_{i}
\frac{\tilde{u}_{i}'}{\tilde{u}_{i}}
+ \sum_{i=1}^{n} \sum_{j=1}^{k_{i}}
\left( \frac{S_{ij}'}{P_{i}^{j}}
- \frac{j S_{ij} P_{i}'}{P_{i}^{j+1}} \right).
\end{equation}

\medskip
\noindent\textbf{The logarithm case.}
Suppose that $t' = a'/a$ (that is, $t = \log(a)$) for a certain
$a \in F$; then $t' \in F$. By Lemma~\ref{lem:degree} one has that for
every monic polynomial $p(t) \in F[t]$ of degree $n > 0$, $(p(t))'$ is
a polynomial in $F[t]$ of degree $n-1$. From this it follows that
$p(t)$ does not divide $(p(t))'$. Applying this reasoning to one of the
$P_{i}$ in the partial fraction decomposition of $v$, one has that in
equation~\eqref{eq:3} a term with $P_{i}^{k_{i}+1}$ appears in the
denominator. Since there is no term that cancels $P_{i}^{k_{i}+1}$,
then it must appear in the expression of $f$, but this is absurd
because $f$ does not depend on $t$. From this it follows that
$v = T \in F[t]$, and therefore
\[
f = T' + \sum_{i=1}^{N} \tilde{c}_{i}
\frac{\tilde{u}_{i}'}{\tilde{u}_{i}}.
\]
Now, reapplying the previous reasoning to $\tilde{u}_{i}$, one has that
there is no term on the right that cancels $\tilde{u}_{i}$, and hence
$\tilde{u}_{i}$ will appear in the expression of $f$. Therefore
$\tilde{u}_{i} \in F$ for every $i = 1, \ldots, N$. We have thus proved
that
\[
T' = f - \sum_{i=1}^{N} \tilde{c}_{i}
\frac{\tilde{u}_{i}'}{\tilde{u}_{i}} \in F.
\]
By Lemma~\ref{lem:degree}, it follows that $T = c t + d$ with
$c \in C_{F}$ and $d \in F$. Finally, then,
\[
f = \sum_{i=1}^{N} \tilde{c}_{i}
\frac{\tilde{u}_{i}'}{\tilde{u}_{i}} + c \frac{a'}{a} + d',
\]
which is an expression of the desired form.

\medskip
\noindent\textbf{The exponential case.}
Suppose now that $t = e^{b} \in F$. Then, again by
Lemma~\ref{lem:degree}, one has that for every monic, irreducible
polynomial $f(t) \in F[t]$ different from $t$, $f(t)$ does not divide
$(f(t))'$. Reasoning in a manner analogous to the case above, one can
conclude that $v \in F[t]$ and $\tilde{u}_{i} \in F$ for every $i$, or
one of them is equal to $t$. One concludes that
\[
v' = f - \sum_{i=1}^{N} \tilde{c}_{i}
\frac{\tilde{u}_{i}'}{\tilde{u}_{i}} \in F,
\]
and, again by Lemma~\ref{lem:degree}, one has that $v \in F$, obtaining
an expression of the desired form.
\end{proof}

Using this theorem one can prove an analogous result in the case in
which the field of constants $C_{F}$ is algebraically closed.

\begin{theorem}\label{thm:liouville-closed}
Let $(F, ')$ be a differential field with field of constants $C_{F}$
algebraically closed and let $f \in F$. If there exists an elementary
extension $E$ of $F$ and $g \in E$ such that $g' = f$, then there exist
$v \in F$, $u_{1}, \ldots u_{n} \in F^{*}$ and
$c_{1}, \ldots, c_{n} \in C_{F}$ such that
\[
f = v' + \sum_{i=1}^{n} c_{i} \frac{u_{i}'}{u_{i}}.
\]
\end{theorem}

The following theorem, which generalises Liouville's result by removing
the restriction on the subfield of constants, was proved by Risch in
1970.

\begin{theorem}[Risch, 1970]\label{thm:risch}
Let $(F, ')$ be a differential field, $C$ its field of constants and
$f \in F$. If there exists an elementary extension $E$ of $F$ and
$g \in E$ such that $g' = f$, then there exist $v \in F$,
$c_{1}, \ldots, c_{n} \in \overline{C}$ and
$u_{1}, \ldots, u_{n} \in F(c_{1}, \ldots, c_{n})^{*}$ such that
\[
f = v' + \sum_{i=1}^{n} c_{i} \frac{u_{i}'}{u_{i}}.
\]
\end{theorem}

A proof of these last two results can be found in~\cite{Bronstein}.

% ============================================================
\section{Examples of functions that are not elementarily integrable}
% ============================================================
In this chapter we will examine the differential field $\C(x)$ with the
standard derivation, which we will denote by the classical symbol $'$.
Moreover, from now on we will consider functions $f$ of a real variable
with complex values. Clearly an $f$ of this type can be decomposed as
\[
f = \Real f + i \Imag f.
\]
The reason for plunging into the complex world is simple: in this world
the trigonometric functions can be traced back to exponential ones,
just as their inverses reduce to logarithms. In particular, the
well-known formula of Euler gives us the following relations:
\[
\sin(x) = \frac{e^{ix} - e^{-ix}}{2i}; \quad
\cos(x) = \frac{e^{ix} + e^{-ix}}{2}; \quad
\tan(x) = \frac{e^{ix} - e^{-ix}}{e^{ix} + e^{-ix}};
\]
\[
\arcsin(x) = \frac{\pi}{2} + i \ln\bigl(ix + \sqrt{1 - x^{2}}\bigr);
\quad
\arctan(x) = \frac{i}{2}\bigl[\ln(1 - ix) - \ln(1 + ix)\bigr].
\]

\begin{theorem}\label{thm:exptranscend}
Let $g \in \C(x)$ be a non-constant rational function. Then the
function $e^{g}$ is transcendental over $\C(x)$.
\end{theorem}

A proof can be found for example in~\cite{DeLellis}.

We now want to derive a criterion to establish in which cases
$f(x) e^{g(x)}$ admits an elementary antiderivative, with $f(x)$ and
$g(x)$ rational functions, $f(x) \neq 0$ and $g(x)$ non-constant. We
write $e^{g} = t$, in this way we have $t'/t = g'$, and we imagine
working in the differential field $\C(x, t)$, a transcendental (but
still elementary) extension of $\C(x)$.

\begin{theorem}\label{thm:criterion}
Let $f, g \in \C(x)$, $f \neq 0$, $g$ non-constant. Then the function
$h(x) = f(x) e^{g(x)}$ admits an elementary antiderivative if and only
if there exists $a \in \C(x)$ such that $f = a' + a g'$.
\end{theorem}

\begin{proof}
($\Leftarrow$) Writing $e^{g} = t$, we have that $t' = t g'$. Suppose
now that there exists $a \in \C(x)$ such that $f = a' + a g'$. Then
\[
(a t)' = a' t + a t' = a' t + a g' t = (a' + a g') t = f t,
\]
and hence $a t$ is an antiderivative of $h$ and is an elementary
function (which lies in $\C(x, t)$, which is elementary by
Definition~\ref{def:elemext}).

($\Rightarrow$) Suppose that $h$ admits an elementary antiderivative
$y \in E$ with $E$ an elementary extension of $\C(x, e^{g})$. Then by
Liouville's theorem there exist $c_{1}, \ldots, c_{n} \in \C$,
$u_{1}, \ldots, u_{n}, v \in \C(x, e^{g})$ such that
\begin{equation}\label{eq:4}
f e^{g} = v' + \sum_{i=1}^{n} c_{i} \frac{u_{i}'}{u_{i}}.
\end{equation}
Since $e^{g}$ is a non-algebraic and exponential element over $\C(x)$,
recalling the proof of Liouville's theorem, one can rewrite
equation~\eqref{eq:4} as
\[
f e^{g} = v' + \sum_{i=1}^{N} d_{i} \frac{v_{i}'}{v_{i}},
\]
with $d_{i} \in \C$, $v_{i} \in \C(x)$ (or one of them equal to
$e^{g}$) and $v \in \C(x)[e^{g}]$. Setting
$d = \sum_{i=1}^{N} d_{i} \frac{v_{i}'}{v_{i}} \in \C(x)$, one has
\[
f e^{g} = d + v'.
\]
Since the left-hand member is a polynomial in $\C(x)[e^{g}]$ of degree
$1$, then the right-hand member must be so too. Hence one has that
$v = s_{0} + s_{1} e^{g}$ with $s_{0}, s_{1} \in \C(x)$. Therefore one
obtains
\[
f e^{g} = d + (s_{0} + s_{1} e^{g})'
= d + s_{0}' + s_{1}' e^{g} + s_{1} g' e^{g},
\]
from which
\begin{equation}\label{eq:5}
e^{g}\bigl(f - s_{1}' - s_{1} g'\bigr) = d + s_{0}'.
\end{equation}
Since $f - s_{1}' - s_{1} g'$ and $d + s_{0}'$ are rational functions
while $e^{g}$ is not, one has that equation~\eqref{eq:5} is satisfied
if and only if
\[
f - s_{1}' + s_{1} g' = 0.
\]
Setting $a := s_{1}$ one has the thesis.
\end{proof}

\subsection*{Notable examples}
Using Theorem~\ref{thm:criterion} it is possible to prove that there
does not exist an elementary antiderivative for some well-known
functions.

\subsubsection*{$e^{x^{2}}$}
Our objective is to deny the existence of a rational function
$a = \frac{P}{Q}$ such that $a' + a g' = e^{x^{2}}$, with $P$ and $Q$
polynomials with complex coefficients, which we assume to be coprime.
Since $g(x) = x^{2}$ and $f(x) = 1$, the identity that they would solve
would be
\[
\frac{P'Q - PQ'}{Q^{2}} + \frac{2xP}{Q} = 1,
\]
which is equivalent to
\begin{equation}\label{eq:6}
P'Q - Q'P + 2xPQ = Q^{2}.
\end{equation}
From this relation one notes that, since $Q$ divides $Q^{2}$, then $Q$
should divide the polynomial $Q'P$. Since however $Q$ and $P$ are
coprime, $Q$ should divide $Q'$. On the other hand this would imply
that $\deg(Q) \le \deg(Q')$, and this is possible only if $Q$ is a
constant polynomial. Hence the identity~\eqref{eq:6} would become
\[
P'(x) + 2xP(x) = 1.
\]
Clearly, however, no polynomial can satisfy this last equality. It is
not difficult to see that nothing changes if instead of $e^{x^{2}}$ we
take $e^{-x^{2}}$, that is, the Gaussian.

\subsubsection*{$\dfrac{e^{x}}{x}$}
Suppose there exists a rational function $a \in \C(x)$ such that
\begin{equation}\label{eq:7}
\frac{1}{x} = a' + a.
\end{equation}
Then, writing $a$ as $\frac{P}{Q}$ with $P, Q \in \C[x]$ coprime, one
can rewrite~\eqref{eq:7} as
\[
\frac{1}{x} = \frac{P'Q - PQ'}{Q^{2}} + \frac{P}{Q}
= \frac{PQ + P'Q - PQ'}{Q^{2}},
\]
which becomes
\[
1 = \frac{x(PQ + P'Q - PQ')}{Q^{2}}
\quad\text{that is,}\quad
Q^{2} = x(PQ + P'Q - PQ').
\]
In particular one has that $Q^{2}$ must divide the numerator, but
$Q^{2}$ does not divide $x$ inasmuch as the first is a product of
second-degree polynomials and the second is a linear polynomial.
Moreover $Q(x)$ divides neither $Q'$ (again because
$\deg(Q) \le \deg(Q')$) nor $P$ (because we assumed that $P, Q$ are
coprime). Hence, as before, one has that $a \in \C[x]$. Equation
\eqref{eq:7} becomes of the form
\[
\frac{1}{x} = P + P'.
\]
Since the right-hand member of the equation is a polynomial and the
left-hand member is a rational function, one can conclude that there
does not exist any $a \in \C[x]$ that satisfies the equality, and hence
by Theorem~\ref{thm:criterion} the function $\frac{e^{x}}{x}$ does not
admit an elementary antiderivative.

\subsubsection*{$\dfrac{\sin x}{x}$}
One finds
\[
\int \frac{\sin x}{x}\, dx
= \int \frac{e^{ix} - e^{-ix}}{2ix}\, dx
= \frac{1}{2i}\left( \int \frac{e^{ix}}{x}\, dx
- \int \frac{e^{-ix}}{x}\, dx \right),
\]
which, by what was seen previously, is not elementarily integrable.

\subsubsection*{$\dfrac{e^{x}}{x^{n}}$, $n \in \N$}
For $n > 0$ we define $I_{n}$ as follows:
\[
I_{n} = \int \frac{e^{x}}{x^{n}}\, dx.
\]
By means of integration by parts we can write
\[
I_{n} = \frac{1}{n-1}\left( \int \frac{e^{x}}{x^{n-1}}\, dx
- \frac{e^{x}}{x^{n-1}} \right)
= \frac{1}{n-1}\left( I_{n-1} - \frac{e^{x}}{x^{n-1}} \right).
\]
One notes that $I_{n}$ is elementarily integrable if and only if
$I_{n-1}$ is. Iterating this procedure $n-1$ times we arrive at proving
that $I_{n}$ is elementarily integrable if and only if $I_{1}$ is, that
is, $e^{x}/x$, which we have proved not to be elementarily integrable.

\subsubsection*{$\dfrac{1}{\ln^{n}(x)}$, $n \in \N$}
As regards this function, with a substitution one has that
\[
\int \frac{dt}{\ln^{n}(t)} \underset{\ln(t)=x}{=}
\int \frac{e^{x}}{x^{n}}\, dx.
\]
And hence, by what was proved previously, $\frac{1}{\ln^{n}(x)}$ too is
not elementarily integrable.

\subsubsection*{Elliptic Integrals}
The theory of elliptic functions arises from the computation of the
length of the ellipse. They can be obtained by inverting functions of
the form $\int dx/\sqrt{P(x)}$ for particular polynomials $P$ with
$P \in \R[x]$, without multiple roots. In general, if $P \in \R[x]$ has
degree $\ge 3$ and has all simple roots, then $\int dx/\sqrt{P(x)}$ is
not elementarily integrable. Since $F := \C(x, \sqrt{P})$ is a
differential field, by Liouville's theorem it is sufficient to prove
that there does not exist an identity of the form
\[
\frac{1}{\sqrt{P(x)}} = v' + \sum_{i=1}^{n} c_{i} \frac{u_{i}'}{u_{i}},
\]
with $c_{1}, \ldots, c_{n} \in \C$ and $u_{1}, \ldots, u_{n}, v \in F$.
From the theory of compact Riemann surfaces one deduces the
impossibility of such an identity. For the details one can
see~\cite{Conrad}.

% ============================================================
\begin{thebibliography}{9}
% ============================================================
\bibitem{Bronstein}
Manuel Bronstein. \emph{Symbolic Integration I: Trascendental
Functions}. Germany: Springer, 2004.

\bibitem{Casolo}
Carlo Casolo. \emph{Dispense del corso di Algebra II}. 2015. URL:
\url{http://web.math.unifi.it/users/casolo/dispense/algebra2_2015.pdf}.

\bibitem{Conrad}
Brian Conrad. \emph{Impossibility theorems for elementary integration}.
1991. URL: \url{http://math.stanford.edu/~conrad/papers/elemint.pdf}.

\bibitem{DeLellis}
Camillo De Lellis. \emph{Il Teorema di Liouville ovvero perch\'e ``non
esiste'' la primitiva di $e^{x^{2}}$}. 2012. URL:
\url{http://www.math.uzh.ch/fileadmin/user/delellis/publikation/Liouville_finale.pdf}.
\end{thebibliography}

\end{document}
